%! AP Statistics — Unit 9, Lesson 9.3 — Guided Notes (Answer Key)
\documentclass[10pt]{article}
\usepackage{apstats-article}
\usepackage{apstats-key}

\begin{document}

\pageheader{Unit 9, Lesson 9.3}{Guided Notes --- Answer Key}
\namedateperiod

\begin{objectivebox}
By the end of this lesson, I will be able to:
\begin{itemize}
    \item[$\square$] \ansline{Interpret a confidence interval for the slope of a regression model in context.}
    \item[$\square$] \ansline{Use a CI for slope to justify or refute a claim about the true slope $\beta$.}
    \item[$\square$] \ansline{Explain how increasing sample size affects the width of a CI for slope.}
\end{itemize}
\end{objectivebox}

\begin{vocabbox}
\termblanklong{CI for slope}
\ansline{An interval estimate of the true slope $\beta$; has the form $b \pm t^* \cdot SE_b$ (UNC-4.AG.1--2).}

\termblanklong{Interpretation of CI for slope}
\ansline{Must include: confidence level (C\%), both endpoints with correct units, and a reference to the \emph{population} (not just the sample).}

\termblanklong{Justifying a claim with a CI}
\ansline{If the claimed value is outside the CI $\Rightarrow$ evidence against that claim; if inside $\Rightarrow$ consistent with the data (UNC-4.AH.1).}

\termblanklong{Width of CI for slope}
\ansline{Width $= 2t^* \cdot SE_b$; larger $n$ decreases $SE_b$, producing a narrower interval when all other things remain the same (UNC-4.AI.1).}
\end{vocabbox}

\begin{hookbox}
Can you use the interval $(0.14,\ 0.52)$ to evaluate the claim $\beta = 0$? What about $\beta = 0.5$?

\ansline{Since 0 is not in $(0.14,\ 0.52)$, this provides evidence against $\beta = 0$. Since 0.5 is in the interval, $\beta = 0.5$ is consistent with the data.}
\end{hookbox}

\begin{notesbox}{Interpreting a CI for Slope (UNC-4.AG)}

\textbf{Interpretation Template:}

``We are \blank{2cm}\% confident that the true slope of the population regression line
relating \blank{4cm} to \blank{4cm} for \blank{4cm} is between \blank{2cm} and
\blank{2cm} \blank{3cm}.''

\vspace{0.1in}
\textbf{Required elements for full AP credit:}
\begin{enumerate}
    \item \ans{Confidence level (C\%)}
    \item \ans{Both endpoints with correct units}
    \item \ans{Reference to the \emph{population} (not just the sample)}
\end{enumerate}

\vspace{0.1in}
\textbf{Worked Example.} A 95\% CI for the slope relating exam score to daily study time
for high school students is $(0.14,\ 0.52)$ points per hour.

\textit{Complete interpretation:}
\ansline{We are 95\% confident that the true slope of the population regression line
relating exam score to daily study time for high school students is between 0.14 and
0.52 points per hour of study.}

\end{notesbox}

\begin{notesbox}{Justifying a Claim Using a CI (UNC-4.AH)}

\textbf{Logic:}
\begin{itemize}
    \item If the claimed value $\beta_0$ is \textbf{NOT} in the CI $\Rightarrow$
          \ans{evidence against that claim (the data are inconsistent with $\beta_0$).}
    \item If the claimed value $\beta_0$ \textbf{IS} in the CI $\Rightarrow$
          \ans{the claim is consistent with the data (plausible).}
\end{itemize}

\vspace{0.1in}
\textbf{Special case --- testing whether $\beta = 0$:}
\begin{itemize}
    \item If 0 is \textbf{not} in the CI $\Rightarrow$
          \ans{evidence that $\beta \ne 0$; a linear relationship exists.}
    \item If 0 \textbf{is} in the CI $\Rightarrow$
          \ans{consistent with no linear relationship; cannot conclude $\beta \ne 0$.}
\end{itemize}

\vspace{0.1in}
\textbf{Practice Table.} 95\% CI for slope $= (0.14,\ 0.52)$

\begin{center}
\renewcommand{\arraystretch}{1.6}
\begin{tabular}{p{3cm}p{3.5cm}p{5.5cm}}
    \textbf{Claimed value} & \textbf{In CI? (Yes/No)} & \textbf{What can we conclude?} \\
    \hline
    $\beta = 0$   & \ans{No}  & \ansline{Evidence that $\beta \ne 0$; a positive linear relationship exists.} \\
    $\beta = 0.5$ & \ans{Yes} & \ansline{Consistent with the data; $\beta = 0.5$ is plausible.} \\
    $\beta = 1.0$ & \ans{No}  & \ansline{Evidence against $\beta = 1.0$; this claim is not supported.} \\
\end{tabular}
\end{center}

\end{notesbox}

\begin{notesbox}{Effect of Sample Size on Width (UNC-4.AI)}

\textbf{Fill in:}

``When all other things remain the same, as sample size \ans{increases}, the standard error
$SE_b$ \ans{decreases}, and the confidence interval for slope becomes \ans{narrower}.''

\vspace{0.05in}
\textbf{Why?}
\ansline{Larger samples contain more information about the population, which reduces variability in the estimate $b$. A smaller $SE_b$ means a smaller margin of error $t^* \cdot SE_b$, producing a narrower interval.}

\end{notesbox}

\end{document}
